\chapter{Implementazione della Pipeline}
\label{chap:implementazione}

Questo capitolo dettaglia le componenti chiave della pipeline di addestramento, con focus sulle funzioni di loss, la logica di training e il bug critico nell'Attention Transfer.

\section{Architettura della Pipeline}

La pipeline supporta quattro modalità operative, gestite dal parametro \texttt{mode} nella configurazione YAML:

\begin{itemize}
    \item \texttt{teacher\_finetune}: fine-tuning del ResNet-50 pre-addestrato su UCF-101.
    \item \texttt{baseline}: addestramento del MobileNet3D da zero con sola cross-entropy.
    \item \texttt{distillation}: KD logit-based con soft targets del teacher.
    \item \texttt{distillation\_at}: KD + Attention Transfer spaziale/temporale.
\end{itemize}

Tutte le modalità condividono la stessa architettura di training: ottimizzatore AdamW~\cite{loshchilov2019adamw}, scheduler cosine annealing, mixed precision con AMP (Automatic Mixed Precision), e gradient clipping a norma 1.0.

\section{Funzione di Loss Combinata}

La loss per la modalità \texttt{distillation\_at} combina tre componenti come descritto dall'Equazione~\ref{eq:total_loss}. Lo pseudo-codice seguente illustra il calcolo:

\begin{algorithm}[H]
\caption{Calcolo della loss combinata KD + AT}\label{alg:loss}
\begin{algorithmic}[1]
\Require Logit student $z^S$, logit teacher $z^T$, label $y$
\Require Feature teacher $\{F^T_i\}_{i=1}^{L}$, feature student $\{F^S_i\}_{i=1}^{L}$
\Require Temperatura $T$, $\alpha$, $\beta_s$, $\beta_t$
\State $p^S \gets \text{softmax}(z^S / T)$, $p^T \gets \text{softmax}(z^T / T)$
\State $\mathcal{L}_{\text{KD}} \gets \alpha \cdot T^2 \cdot D_{\text{KL}}(p^S \| p^T) + (1 - \alpha) \cdot \text{CE}(z^S, y)$
\State $\mathcal{L}_{\text{AT}} \gets 0$
\For{$i = 1$ \textbf{to} $L$}
    \State $A^T_s \gets \text{L2\_norm}(\text{mean}_{c,t}((F^T_i)^2))$ \Comment{Mappa spaziale teacher}
    \State $A^S_s \gets \text{L2\_norm}(\text{mean}_{c,t}((F^S_i)^2))$ \Comment{Mappa spaziale student}
    \State $A^S_s \gets \text{interp2d}(A^S_s, \text{size}=A^T_s.\text{shape})$ \Comment{Allineamento risoluzione}
    \State $\mathcal{L}_{\text{AT}} \gets \mathcal{L}_{\text{AT}} + \beta_s \cdot \text{MSE}(A^T_s, A^S_s)$
    \State $A^T_t \gets \text{L2\_norm}(\text{mean}_{c,h,w}((F^T_i)^2))$ \Comment{Profilo temporale teacher}
    \State $A^S_t \gets \text{L2\_norm}(\text{mean}_{c,h,w}((F^S_i)^2))$ \Comment{Profilo temporale student}
    \State $\mathcal{L}_{\text{AT}} \gets \mathcal{L}_{\text{AT}} + \beta_t \cdot \text{MSE}(A^T_t, A^S_t)$
\EndFor
\State \Return $\mathcal{L}_{\text{KD}} + \mathcal{L}_{\text{AT}}$
\end{algorithmic}
\end{algorithm}

\section{Il Bug del Block 5 (``Key~5'')}
\label{sec:at_bug}

Un bug critico ha impattato significativamente le prime iterazioni dell'Attention Transfer. In una versione iniziale del codice, i \texttt{FEATURE\_BLOCKS} del teacher erano configurati come $[2, 3, 4, 5]$, dove il blocco 5 corrisponde alla \textbf{testa di classificazione} (\texttt{ResNetBasicHead}), non ad uno stage residuale. Questo blocco restituisce un tensore mono-dimensionale (\texttt{[B, 101]}) dopo il pooling e la proiezione lineare.

Le conseguenze di questo errore erano molteplici:
\begin{itemize}
    \item Il tensore del blocco 5, non essendo una feature map volumetrica, produceva mappe di attenzione costanti (tutti i valori a 1.0 dopo la normalizzazione).
    \item La loss AT su questa coppia era identicamente zero, rendendo il segnale AT di fatto nullo.
    \item L'errore era \textbf{silenzioso}: nessuna eccezione, nessun warning, ma una convergenza dell'AT identica al training senza AT.
\end{itemize}

La correzione ha rimosso il blocco 5 dall'elenco, settando \texttt{FEATURE\_BLOCKS = [2, 3, 4]}.

\section{Gestione dell'Interpolazione Spaziale}

Poiché le risoluzioni spaziali di teacher e student differiscono (Tabella~\ref{tab:feature_pairing}), la mappa di attenzione spaziale dello student viene ridimensionata per corrispondere a quella del teacher. L'interpolazione avviene in 2D (non 3D), coerentemente con il fatto che le mappe spaziali hanno già aggregato la dimensione temporale:

\begin{enumerate}
    \item Si fa il reshape della mappa $A^S \in \mathbb{R}^{B \times H_S W_S}$ a $[B, 1, H_S, W_S]$.
    \item Si applica \texttt{F.interpolate} con mode \texttt{bilinear} e \texttt{align\_corners=False}.
    \item Si riduce nuovamente alla dimensione della mappa teacher $[B, H_T \cdot W_T]$.
\end{enumerate}

Per la dimensione temporale, non è necessaria alcuna interpolazione poiché $T$ è identica tra teacher e student a tutti i livelli.

\section{Warmup per la Knowledge Distillation}

È stato identificato che le prime epoche di addestramento con KD beneficiano significativamente di una fase di ``warmup'' in cui la loss KD viene esclusa e lo student viene addestrato con sola cross-entropy. Questo permette ai parametri del modello di raggiungere una regione ragionevole dello spazio dei pesi prima di essere esposti al segnale di distillazione, evitando instabilità numeriche e convergenza prematura in minimi locali sfavorevoli. Nel progetto, un warmup di 10 epoche è stato adottato come default per tutti gli esperimenti di distillazione.

\section{Mixed Precision Training}

La pipeline utilizza AMP (\textit{Automatic Mixed Precision}) per accelerare il training e ridurre l'occupazione di memoria GPU. Un problema specifico riscontrato è stato l'overflow in FP16 durante il calcolo delle mappe di attenzione: l'operazione $F^2$ su feature map con valori grandi genera overflow nel range di float16. La soluzione adottata è il casting esplicito a \texttt{float32} prima dell'operazione di potenza:

\begin{verbatim}
# Cast esplicito per stabilità numerica con AMP
f = f.float()  # float32
return (f * f).mean(dim=...)
\end{verbatim}

Questo garantisce stabilità numerica senza disabilitare AMP per l'intero training.
